%! AP Statistics — Unit 9, Lesson 9.2 — Homework
\documentclass[10pt]{article}
\usepackage{apstats-article}
\usepackage{apstats-boxes}
\usepackage{grading}
\gradingassignment{apstats-u09-l02-hw}
\gradingcourse{advanced_statistics}

\begin{document}

\pageheader{Unit 9, Lesson 9.2}{Homework}

\vspace{0.1in}

\textbf{Problem 1: Sleep and Reaction Time}

\vspace{0.05in}
A random sample of 22 college students was used to study the relationship between hours
of sleep per night ($x$) and reaction time in milliseconds ($y$). Computer output is
shown below.

\vspace{0.1in}
\begin{center}
\renewcommand{\arraystretch}{1.3}
\begin{tabular}{lrrrr}
    \textbf{Predictor} & \textbf{Coef} & \textbf{SE Coef} & \textbf{T} & \textbf{P} \\
    \hline
    Constant & 412.3 & 28.4 & 14.52 & 0.000 \\
    Sleep    & $-18.6$ & 5.3 & $-3.51$ & 0.002 \\
    \hline
    \multicolumn{5}{l}{$s = 24.1$ \quad $n = 22$}
\end{tabular}
\end{center}

\vspace{0.1in}
\begin{enumerate}[label=\alph*.]

    \boxguard[15]
    \item Identify $b$ and $SE_b$ from the output.
          \gradedhide{\writelines{1}\gradedright{$b = $~\gradedblank[id=hw1a1, type=short, points=1, answer=-18.6, check=numeric, tolerance=0]{0.75in}{0.2in}\quad $SE_b = $~\gradedblank[id=hw1a2, type=short, points=1, answer=5.3, check=numeric, tolerance=0]{0.75in}{0.2in}}}

    \boxguard[15]
    \item For this problem, assume: the data come from a random sample; the residual plot
          shows random scatter with roughly constant spread; and $n = 22$ is large enough
          for normality (stated). Briefly note how each LINER condition is addressed.
          \gradedhide{\writeline\\\gradedhere{\gradedblank[id=hw1b, type=rubric, points=3, display={\textbf{L}: Random scatter in residual plot --- linear condition met. \textbf{I}: Random sample of college students; $n = 22 \le 10\%$ of all students --- independent. \textbf{N}: $n = 22$; problem states conditions are met --- normal condition assumed. \textbf{E}: Roughly constant spread in residual plot --- equal variance met. \textbf{R}: Random sample --- random condition met.}]{\linewidth+6pt}{\gradedlinesheight{4}}}\writelines{9}}

    \boxguard[15]
    \item Construct a 95\% CI for $\beta$ using $t^* = 2.086$ ($df = 20$). Show all work.
          \gradedhide{\writelines{3}\gradedright{(~\gradedblank[id=hw1c1, type=short, points=1, answer=-29.7, check=numeric, tolerance=0.1]{0.8in}{0.2in},~\gradedblank[id=hw1c2, type=short, points=1, answer=-7.5, check=numeric, tolerance=0.1]{0.8in}{0.2in}~)}}

    \boxguard[15]
    \item Interpret the CI in context.
          \gradedhide{\writeline\\\gradedhere{\gradedblank[id=hw1d, type=rubric, points=2, display={95\% confident the true slope (reaction time vs.\ hours of sleep, college students) is between $-29.7$ and $-7.5$ ms per hour.}]{\linewidth+6pt}{\gradedlinesheight{1}}}\writelines{9}}

    \boxguard[15]
    \item Does the CI provide evidence of a negative linear relationship between hours of
          sleep and reaction time? Explain.
          \gradedhide{\writeline\\\gradedhere{\gradedblank[id=hw1e, type=rubric, points=2, display={\bfseries Yes: the whole CI is below 0 (0 not in it), so $\beta < 0$: more sleep, faster reactions.}]{\linewidth+6pt}{\gradedlinesheight{1}}}\writelines{9}}

\end{enumerate}

\vspace{0.2in}

\textbf{Problem 2: Temperature and Heating Fuel Cost}

\vspace{0.05in}
A researcher studies 18 cities, recording average daily high temperature in January
($x$, ${}^\circ$F) and monthly heating fuel cost ($y$, \$/month).
Output gives $b = -4.8$, $SE_b = 1.1$, and $n = 18$.
Use $t^* = 2.120$ ($df = 16$). Assume all LINER conditions are met.

\vspace{0.1in}
\begin{enumerate}[label=\alph*.]

    \boxguard[15]
    \item Construct a 95\% CI for the true slope $\beta$. Show all work, then interpret
          the interval in context.
          \gradedhide{\writeline\\\writeline\gradedright{(~\gradedblank[id=hw2a1, type=short, points=1, answer=-7.13, check=numeric, tolerance=0.02]{0.8in}{0.2in},~\gradedblank[id=hw2a2, type=short, points=1, answer=-2.47, check=numeric, tolerance=0.02]{0.8in}{0.2in}~)}\\\gradedhere{\gradedblank[id=hw2a3, type=rubric, points=2, display={We are 95\% confident that the true slope of the population regression line relating monthly heating fuel cost to January temperature for all cities is between $-\$7.13$ and $-\$2.47$ per ${}^\circ$F.}]{\linewidth+6pt}{\gradedlinesheight{2}}}\writelines{9}}

    \boxguard[15]
    \item A utility company claims $\beta = -5$ (each 1${}^\circ$F warmer costs \$5 less
          per month). Is this claim consistent with your CI? Explain.
          \gradedhide{\writeline\\\gradedhere{\gradedblank[id=hw2b, type=rubric, points=2, display={Yes. The value $\beta = -5$ falls inside the CI $(-7.13,\ -2.47)$, so the utility company's claim is consistent with the data.}]{\linewidth+6pt}{\gradedlinesheight{1}}}\writelines{9}}

    \boxguard[15]
    \item If the study had used 72 cities instead of 18, would you expect the CI to be
          narrower or wider? Explain using the formula for $SE_b$.
          \gradedhide{\writeline\\\gradedhere{\gradedblank[id=hw2c, type=rubric, points=2, display={\bfseries Narrower: $SE_b = s/(s_x\sqrt{n-1})$ is smaller for larger $n$, so a smaller margin of error.}]{\linewidth+6pt}{\gradedlinesheight{1}}}\writelines{9}}

\end{enumerate}

\end{document}
